% sec-05: chronology, simulations, RASolute 302.  Figures 4, 7, 16.
\section{Trial Evidence}

\subsection{The Chronology}

\begin{appfloat}[!tb]
\begin{appfig}[mermaid-type / gantt]
% Time axis of 15 months at 0.78cm per month.  Six author rows 0.62 apart, then
% a rule, then the external row.  The two 2025 overlaps are the only bars in the
% darker fill, because the overlap is the figure's point.
\foreach \m/\lab in {0/{Jun 25}, 3/{Sep 25}, 6/{Dec 25}, 9/{Mar 26}, 12/{Jun 26}, 14/{Aug 26}}{
  \draw[pagraym,line width=0.3pt] (\m*0.78,0.42) -- (\m*0.78,-4.35);
  \node[font=\tiny\sffamily,text=protogray,anchor=south] at (\m*0.78,0.45) {\lab};}
\fill[pagrayl] (0.78,-0.32) rectangle (2.34,-1.56);
\node[mmbar,minimum width=1.56cm,anchor=west] at (0,0) {};
\node[font=\tiny\sffamily,anchor=west] at (1.64,0) {Drug identification, 40 PDAC meta-analyses};
\node[mmbark,minimum width=1.56cm,anchor=west] at (0.78,-0.62) {};
\node[font=\tiny\sffamily,anchor=west] at (2.42,-0.62) {QSP simulation, 10 arms, 250 ODEs};
\node[mmbark,minimum width=1.56cm,anchor=west] at (2.34,-1.24) {};
\node[font=\tiny\sffamily,anchor=west] at (3.98,-1.24) {Digital twin, 1000 patients, 55 tests};
\node[mmbar,minimum width=2.34cm,anchor=west] at (7.80,-1.86) {};
\node[font=\tiny\sffamily,anchor=west] at (10.22,-1.86) {Protocols, IND, bill drafts};
\node[mmbar,minimum width=0.78cm,anchor=west] at (10.14,-2.48) {};
\node[font=\tiny\sffamily,anchor=west] at (11.00,-2.48) {Two funding applications, 5 days apart};
\node[mmbar,minimum width=0.78cm,anchor=west] at (10.92,-3.10) {};
\node[font=\tiny\sffamily,anchor=west] at (11.78,-3.10) {Ten application file sets};
\draw[pagrayd,line width=0.5pt] (0,-3.55) -- (13.0,-3.55);
\node[mmbarg,minimum width=0.30cm,anchor=west] at (8.58,-3.95) {};
\node[font=\tiny\sffamily,anchor=west] at (8.96,-3.95) {RASolute 302 readout, independent};
\node[mmlanetitle,anchor=east,align=right] at (-0.15,-1.24) {Author\\work};
\node[mmlanetitle,anchor=east,align=right] at (-0.15,-3.95) {External};
\node[font=\tiny\itshape,text=protogray,anchor=north west,text width=112mm]
  at (0,-4.55) {The shaded band marks the two months in which two activities ran
  at once. Bar length is elapsed time, not effort.};
\end{appfig}
\vspace{-0.7cm}
\figcaption{Fourteen months of author work against the one external readout. Bar
length is elapsed time, not effort, and the two 2025 overlaps are what a single
operator could not have run sequentially.}
\end{appfloat}

In June 2025 deep-research meta-analysis over forty pancreatic cancer literature
areas, more than 400,000 words, produced a ranked shortlist from which
daraxonrasib emerged as the leading candidate \cite{daraxident}. In August 2025
a quantitative systems pharmacology simulation, ten arms and 250 ordinary
differential equations per patient, returned median overall survival of 12.8
months against 5.4 for chemotherapy \cite{qspmetpancre}. In October 2025 a
1000-patient digital twin returned a best progression-free survival hazard ratio
of 0.31 with an ASME V\&V 40 and ICH M15 aligned credibility score of 81.9
across 55 verification notebook tests \cite{daraxsims,asmevv40,ichm15}.

In May 2026, nine months after the QSP run, RASolute 302 reported 13.2 months
against 6.6 in the RAS G12 population of previously treated metastatic
pancreatic cancer \cite{rasolute302}. Between June and July 2026 the regulatory
package was written and two complete NIH funding applications were deposited
five days apart \cite{fundapp1,fundapp2}.

\subsection{What the Numbers Are, and What They Are Not}

\begin{apptable}
\begin{tabularx}{\textwidth}{>{\raggedright\arraybackslash}p{2.6cm} >{\raggedright\arraybackslash}p{1.5cm} >{\raggedright\arraybackslash}p{1.6cm} >{\raggedright\arraybackslash}p{1.5cm} >{\raggedright\arraybackslash}X}
\headrow \thead{Source} & \thead{Date} & \thead{Test arm} & \thead{Control} & \thead{Stated limitation, in the authors' own words}\\
\midrule
Empirical triplicate, 100,000 patients & Jul 2025 & G3+ AE 8.0\% & G3+ AE 25.0\% & KRAS G12C log file favours experimental while the KRAS-mutant report favours control; trial-to-trial variability \\
QSP simulation, 10 arms & Aug 2025 & 12.8 mo mOS, HR 0.25 & 5.4 mo mOS & Assumes no acquired resistance and ideal pharmacodynamics; G3+ AE consistently high in all arms \\
Digital twin, 1000 patients & Oct 2025 & 12.1 mo mOS, mPFS HR 0.31 & not applicable & No patient-specific PK, PD, or tumour growth parameters; simple Emax model, no immune compartments \\
RASolute 302 & May 2026 & 13.2 mo mOS & 6.6 mo mOS & Metastatic and previously treated; silent on the resectable setting \\
\end{tabularx}
\end{apptable}
\tabcap{Four sources with the limitation each set of authors stated, carried in
the same row as the result. No row can be read without the caveat that qualifies
it, which is the condition on citing any of them.}

The 2025 simulated ratio of 2.4-fold and the 2026 observed ratio of 2.0-fold are
close. That proximity is a chronology observation and a hypothesis-supporting
one. It is \textbf{not} a validation claim, and three differences are material:
1000 simulated against 241 enrolled; a combination against a single agent; and
KRAS G12C against a primarily G12D and G12V population. Those clauses appear in
all ten applications, in the same paragraph as the numbers, for the same reason
they appear here.

\begin{appfloat}[!tb]
\begin{appfig}[d2-type / nested containers]
% Four nested containers drawn on the background layer so no container edge
% crosses a leaf.  Titles sit at each container's top-left, outside the child
% field, 1.2mm above the edge.
\node[d2soft,text width=40mm,minimum height=7.5mm] (l1) at (0,0) {QSP, 10 arms, 250 ODEs\\mOS 12.8 vs 5.4};
\node[d2soft,text width=40mm,minimum height=7.5mm,anchor=north west] (l2) at ([yshift=-2.6mm]l1.south west) {Empirical triplicate\\100,000 patients};
\begin{scope}[on background layer]
\node[d2cont2,fit=(l1)(l2),inner sep=6pt] (c1) {};
\end{scope}
\node[d2title2,anchor=south west] at ([yshift=1.2mm]c1.north west) {Tier 1: in silico exploration};
\node[d2soft,text width=40mm,minimum height=7.5mm,anchor=north west] (l3) at ([yshift=-8mm]c1.south west) {Digital twin, credibility 81.9\\55 verification tests};
\node[d2soft,text width=40mm,minimum height=7.5mm,anchor=north west] (l4) at ([yshift=-2.6mm]l3.south west) {ASME V\&V 40 and ICH M15\\alignment};
\begin{scope}[on background layer]
\node[d2cont,fit=(c1)(l3)(l4),inner sep=7pt] (c2) {};
\end{scope}
\node[d2title,anchor=south west] at ([yshift=1.2mm]c2.north west) {Tier 2: verified computation};
\node[d2key,text width=40mm,minimum height=7.5mm,anchor=north west] (l5) at ([yshift=-8mm]c2.south west) {RASolute 302\\mOS 13.2 vs 6.6, RAS G12};
\begin{scope}[on background layer]
\node[d2cont,fit=(c2)(l5),inner sep=7pt] (c3) {};
\end{scope}
\node[d2title,anchor=south west] at ([yshift=1.2mm]c3.north west) {Tier 3: peer-reviewed clinical result};
\node[d2gray,text width=40mm,minimum height=7.5mm,anchor=north west] (l6) at ([yshift=-8mm]c3.south west) {This Phase 1, up to 18\\participants, no result yet};
\begin{scope}[on background layer]
\node[d2ghost,fit=(c3)(l6),inner sep=7pt] (c4) {};
\end{scope}
\node[d2title2,anchor=south west] at ([yshift=1.2mm]c4.north west) {Tier 4: proposed clinical research};
\node[font=\tiny\itshape,text=protogray,anchor=north west,text width=58mm]
  at ([xshift=54mm,yshift=-6mm]c4.north west) {The outer container is dashed
  because it holds no evidence yet. Nothing inside a container may be cited as
  though it were the container: a tier 1 result is not a tier 3 result, and the
  proposed trial is not a result at all.};
\end{appfig}
\vspace{-0.7cm}
\figcaption{Four evidence tiers drawn as containment rather than as a ladder.
Nothing inside a container may be cited as though it were the container, which
is the whole of the paper's evidentiary discipline.}
\end{appfloat}

\subsection{Cost and Schedule}

\begin{apptable}
\begin{tabularx}{\textwidth}{>{\raggedright\arraybackslash}p{4.3cm} >{\raggedright\arraybackslash}p{2.6cm} >{\raggedright\arraybackslash}p{2.5cm} >{\raggedright\arraybackslash}X}
\headrow \thead{Work product} & \thead{Industry cost} & \thead{Industry schedule} & \thead{This author}\\
\midrule
Empirical virtual trial, triplicate & above \$120,000 per run & 4.5 months & \$36,330 estimated, about one month \\
Ten-arm QSP virtual trial & above \$2,000,000 & several months to over a year & \$36,330 estimated, about one month \\
Digital-twin simulator platform & \$28,000 to \$700,000 & 3 to 16 weeks & \$36,330 estimated, about one month \\
\end{tabularx}
\end{apptable}
\tabcap{Three work products against their industry cost and schedule benchmarks.
The author figures are estimates from the source set and are reported as such,
not as audited accounts.}

\subsection{Provenance}

\begin{appfloat}[!tb]
\begin{appfig}[graphviz-type / citation graph]
% Two clusters, left at x = 1.2 and right at x = 8.6, so the two lines are
% visibly parallel.  The corridor between them holds only the two dashed
% crossing edges; its emptiness is the argument.
\foreach \i/\y/\lab in {1/0/{VVUQ Bill v1}, 2/-1.3/{Verification Before Generation v2},
  3/-2.6/{H. R. 9510 v3}, 4/-3.9/{H. R. 9510 v4}, 5/-5.2/{H. R. 9510 v5},
  6/-6.5/{From H. R. 9510 to Federal Law}, 7/-7.8/{Earning the Clinician's Trust},
  8/-9.1/{Earning the Congress's Vote}}{
  \node[gvnode,text width=22mm] (b\i) at (1.2,\y) {\lab};}
\begin{scope}[on background layer]
\node[gvcluster,fit=(b1)(b8),inner sep=7pt] (cb) {};
\end{scope}
\node[gvctitle,anchor=south west] at ([yshift=1mm]cb.north west) {Legislative line};
\foreach \i/\y/\lab/\sty in {1/0/{Phase 1 Protocol}/gvkey, 2/-1.3/{Phase 2 Protocol}/gvnode,
  3/-2.6/{PI LLM Adoption Guide}/gvnode, 4/-3.9/{Phase 1 LLM Guidance}/gvnode,
  5/-5.2/{IND application}/gvkey, 6/-6.5/{Funding Application v1}/gvnode,
  7/-7.8/{Funding Application v2}/gvnode}{
  \node[\sty,text width=22mm] (p\i) at (8.6,\y) {\lab};}
\begin{scope}[on background layer]
\node[gvcluster2,fit=(p1)(p7),inner sep=7pt] (cp) {};
\end{scope}
\node[gvctitle,anchor=south west] at ([yshift=1mm]cp.north west) {Trial line};
\foreach \i/\j in {1/2, 2/3, 3/4, 4/5, 5/6, 6/7, 7/8}{\draw[gvedge] (b\i) -- (b\j);}
\draw[gvedge] (p1) -- (p2);
\draw[gvedge] (p1) to[out=0,in=0,looseness=1.1] (p3);
\draw[gvedge] (p3) -- (p4);
\draw[gvedge] (p4) -- (p5);
\draw[gvedge] (p2) to[out=0,in=0,looseness=0.9] (p5);
\draw[gvedge] (p5) -- (p6);
\draw[gvedge] (p6) -- (p7);
\draw[gvedged] (b5.east) -- node[font=\tiny,pos=0.5,fill=protowhite,inner sep=1pt] {verification frame} (p1.west);
\draw[gvedged] (b7.east) -- node[font=\tiny,pos=0.5,fill=protowhite,inner sep=1pt] {clinician objections} (p3.west);
\end{appfig}
\vspace{-0.7cm}
\figcaption{Fourteen deposited works and the citation edges between them. The
bill line and the trial line run in parallel and meet only twice, which is why
the legislative drafts are not on the trial's path.}
\end{appfloat}

Two edges cross the corridor: the verification frame from the five bill versions
into the Phase 1 protocol, and the clinician objections into the principal
investigator adoption guide. Everything else in the legislative line stays
there.

That separation answers a question a reviewer will ask. The trial does not
depend on legislation passing. The bill work supplied a verification vocabulary
and a catalogue of objections, and nothing in the protocol, the investigational
new drug application, or the ten applications is contingent on H. R. 9510 being
enacted \cite{hr9510billv5,clintrust}.
